虽然构造函数（constructor）负责对象创建时的各种情形，但 C++ 也提供了一种处理对象销毁的方式。C++ 不像许多流行的编程语言那样提供任何形式的垃圾回收，但得益于精确的生存期（lifetime）规范，你可以确信对象何时会被销毁。

每个类都有一个被称为析构函数（destructor）的特殊成员函数。如果你不编写它，编译器会准备一个默认实现。当对象的生存期结束时，会调用析构函数。在大多数情况下，这意味着对象离开了作用域（对于栈上分配的变量），或者调用了 \texttt{delete} 运算符（对于堆上分配的变量）。此外，当你有一个用户自定义类时，它会自动为其数据成员调用析构函数。关于生存期的更多信息，请参阅 C++ 参考页\footnote{\url{https://en.cppreference.com/w/cpp/language/destructor}} 上的精彩总结。

\section{基础}

在我们继续之前，应该扩展一下术语。到目前为止，我用“对象（object）”来指代某种类型的实体，并依赖我们的“直觉”来访问这些实体。但 C++ 标准用以下术语定义了一个对象（简化自 C++ 草案 — intro.object\footnote{\url{https://timsong-cpp.github.io/cppwp/n4868/intro.object\#1}}）：

\begin{myQuote}
• C++ 程序中的构造会创建、销毁、引用、访问和操作对象。

• 对象由定义、new 表达式、隐式创建对象的操作创建，或者在创建临时对象时创建。对象在其构造期间、整个生存期内以及销毁期间占据一片存储区域。
\end{myQuote}

继续说：

\begin{myQuote}
• 对象可以有一个名字，

• 对象有存储期（storage duration），它影响其生存期，

• 对象有一个类型，

• 对象可以包含其他对象，称为子对象（subobject）。子对象可以是成员子对象、基类子对象或数组元素。
\end{myQuote}

下面是一个处理对象生存期结束情形的析构函数基本场景：

\begin{codeboxt}[cpp]{代码清单 5.1 一个记录日志的析构函数（在 Compiler Explorer 上运行）}
#include <iostream>
#include <string>
class Product {
public:
explicit Product(const char* name, unsigned id)
: name_(name), id_(id) {
std::cout << name << ", id " << id << '\n';
}
~Product() { std::cout << name_ << " destructor...\n"; }
std::string Name() const { return name_; }
unsigned Id() const { return id_; }
private:
std::string name_;
unsigned id_;
};
\end{codeboxt}

这个例子包含以下特殊成员函数：

\begin{codebox}[cpp]
~Product() { std::cout << name_ << " destructor...\n"; }
\end{codebox}

这种语法很独特，它没有参数，并且带有 \texttt{\textasciitilde} 前缀。一个类中也只能有一个析构函数。此外，析构函数不返回任何值。

现在，让我们创建该类型的两个对象：

\begin{codeboxt}[cpp]{代码清单 5.1 一个记录日志的析构函数（续）（在 Compiler Explorer 上运行）}
int main() {
{
Product tvset("TV Set", 123);
}
{
Product car("Mustang", 999);
}
}
\end{codeboxt}

在我们的例子中，构造函数和析构函数被用来执行日志输出。当你运行这个例子时，会看到以下输出：

\begin{codebox}[cpp]
TV Set, id 123
TV Set destructor...
Mustang, id 999
Mustang destructor...
\end{codebox}

我特意将对象（创建在栈上）放在各自独立的作用域中，这样它们的生存期会在作用域结束时终止。另一方面，如果我们有这样的代码：

\begin{codebox}[cpp]
int main() {
Product tvset("TV Set", 123);
Product car("Mustang", 999);
}
\end{codebox}

那么 \texttt{tvset} 和 \texttt{car} 共享同一个生存期作用域，因此我们预期会看到以下输出：

\begin{codebox}[cpp]
TV Set, id 123
Mustang, id 999
Mustang destructor...
TV Set destructor..
\end{codebox}

如你所见，析构函数按它们被创建的逆序调用。这是因为栈是一个 LIFO 结构（Last In, First Out，后进先出）。\texttt{tvset} 先被创建并压入栈，然后 \texttt{car} 被压入。当函数离开作用域时，栈被清空，按逆序取出元素。所以 \texttt{car} 先被销毁，然后才是 \texttt{tvset}。下面的示意图展示了这一点：从栈中添加和移除对象。

\section{分配在堆上的对象}

在之前的例子中，我使用了创建在栈上的对象。为了更清楚地理解析构函数，最好讨论一下对象位于堆上的情形。在这种情况下，只有在通过 \texttt{delete} 运算符释放内存时，才会调用析构函数。

请看以下代码片段：

\begin{codebox}[cpp]
{
Product* ptr = new Product("TV Set", 123);
}
// !!
\end{codebox}

\texttt{ptr} 是一个指向堆上分配的对象的指针。但我没有调用 \texttt{delete}，因此析构函数不会被调用！此外，我还造成了内存泄漏，因为内存也没有被释放。作用域结束后，\texttt{ptr} 离开了作用域，但由于它是指针，内存仍然存在但无法访问。

为了修复这个问题，我们必须调用 \texttt{delete}（针对单个元素）或 \texttt{delete []}（针对数组）。

\myGraphic{0.9}{images/fig_p091_1.png}{}

\begin{codebox}[cpp]
{
Product* ptr = new Product("TV Set", 123);
// use ptr...
delete ptr;
}
\end{codebox}

类似地，对于数组：

\begin{codebox}[cpp]
{
Product* arr = new Product[10]("TV Set", 123);
// use...
delete [] arr;
}
\end{codebox}

由于很容易忘记正确地释放堆内存，最好使用智能指针（smart pointer）将分配与所有权封装在一起。

\begin{codebox}[cpp]
{
std::unique_ptr<Product> ptr = std::make_unique<Product>("box", 1);
// use ptr...
}
\end{codebox}

现在，当 \texttt{ptr} 离开作用域时，它是“智能”的，知道还必须在所存储的指针上调用 \texttt{delete}。在例子中，我使用了 \texttt{unique\_ptr}，它是 C++ 标准库中智能指针的基本形式。它封装了指向对象的指针并保持对其的“唯一”所有权。如果你想在系统中传递指针并拥有多个“所有者”，那么可以使用 \texttt{shared\_ptr}。

智能指针同样适用于数组版本：

\begin{codebox}[cpp]
{
// create 10 Products
std::unique_ptr<Product[]> ptr = std::make_unique<Product[]>(10);
// use...
}
\end{codebox}

这一次，\texttt{unique\_ptr} 确保调用 \texttt{delete []}。你可以在 Compiler Explorer\footnote{\url{https://godbolt.org/z/Ps9Ye79zc}} 上试验这个例子。

关于智能指针的更多信息，请参阅我的博客系列：6 Ways to Refactor new/delete into \texttt{unique\_ptr} - C++ Stories\footnote{\url{https://www.cppstories.com/2021/refactor-into-uniqueptr/}}，以及更多关于智能指针的文章 @C++Stories\footnote{\url{https://www.cppstories.com/tags/smart-pointers/}}。

\begin{myInfoBox}
此外，Modern C++ 强烈建议避免使用原始的 \texttt{new} 和 \texttt{delete}。得益于许多库容器、包装器和智能指针，几乎不再需要依赖这些底层的例程。参见这条 C++ Core 指南：R.11：避免显式调用 \texttt{new} 和 \texttt{delete}\footnote{\url{https://isocpp.github.io/CppCoreGuidelines/CppCoreGuidelines\#r11-avoid-calling-new-and-delete-explicitly}}。本节中带 \texttt{new} 的代码仅作示意用途。
\end{myInfoBox}

\section{析构函数与数据成员}

如果你有一个类，那么默认情况下，它的析构函数会为所有数据成员调用其析构函数：

\begin{codeboxt}[cpp]{代码清单 5.2 嵌套的析构函数调用（在 Compiler Explorer 上运行）}
#include <iostream>
#include <string>
class Product {
// defined as in the previous example...
};
class Wrapper {
public:
Wrapper() : prod_("internal", 123) { std::cout << "Wrapper()\n"; }
~Wrapper() { std::cout << "~Wrapper()\n"; }
private:
Product prod_;
};
int main() {
Wrapper w;
}
\end{codeboxt}


在例 5.2 中，\texttt{Wrapper} 类将 \texttt{Product} 作为一个数据成员。

输出：

\begin{codebox}[cpp]
internal, id 123
Wrapper()
~Wrapper()
internal destructor...
\end{codebox}

你会注意到，\texttt{internal destructor} 与 \texttt{\textasciitilde Wrapper()} 的调用一起被调用。

另一方面，当你的对象作为某个其他类的指针成员时，它会离开作用域，但它所指向的数据不会。所以你必须注意你的指针数据成员，并在合适的地方调用它们的 \texttt{delete}。看下面的例子：

\begin{codeboxt}[cpp]{代码清单 5.3 指针数据成员与析构函数（在 Compiler Explorer 上运行）}
class Wrapper {
public:
Wrapper() : prod_(new Product("internal", 123)) {
std::cout << "Wrapper()\n";
}
~Wrapper() {
delete prod_;
std::cout << "~Wrapper()\n";
}
private:
Product *prod_;
};
int main() {
Wrapper w;
}
\end{codeboxt}

这一次我必须手动调用 \texttt{delete prod\_} 来正确地释放数据成员。这又是智能指针派上用场的一种情形，因为它会自动销毁底层对象。

\section{虚析构函数与多态}

析构函数还有一个特性，在继承和多态（polymorphism）中扮演着重要角色。

根据 Wikipedia\footnote{\url{https://en.wikipedia.org/wiki/Polymorphism_(computer_science)}}：

\begin{myQuote}
在编程语言理论和类型理论中，多态是为不同类型的实体提供单一接口，或使用单一符号来表示多个不同类型。
\end{myQuote}

在 C++ 中，这个定义意味着如果你有一个指向基类（base class）的指针或引用，当你调用成员函数时，编译器会调用派生类（derived class）中的实现（如果可用）。C++ 通过 \texttt{virtual} 函数实现这种技术。我们可以用下面这段朴素（但简单）的 C++ 代码来演示：

\begin{codeboxt}[cpp]{代码清单 5.4 虚析构函数，基类，不正确（在 Compiler Explorer 上运行）}
class Product {
public:
explicit Product(const char* name) : name_(name) {
std::cout << name << '\n';
}
~Product() { std::cout << name_ << " destructor...\n"; }
std::string Name() const { return name_; }
virtual double CalculateMass() const = 0;
private:
std::string name_;
};
\end{codeboxt}

上面的 \texttt{Product} 类型声明了一个 \texttt{virtual} 成员函数。我们可以声明派生类，然后提供该虚成员函数的实现。这让编译器能够根据类型调用正确的函数，并赋予这种“多态”语义。

请看：

\begin{codeboxt}[cpp]{代码清单 5.4 虚析构函数，派生类（在 Compiler Explorer 上运行）}
struct BoxProduct : public Product {
using Product::Product; // inheriting ctor
~BoxProduct() { std::cout << "~BoxProduct...\n"; }
double CalculateMass() const override { return 10.0; }
};
struct FluidProduct : public Product {
using Product::Product; // inheriting ctor
~FluidProduct() { std::cout << "~FluidProduct...\n"; }
double CalculateMass() const override { return 100.0; }
};
\end{codeboxt}

\texttt{CalculateMass} 函数有两个独立而简单的实现\footnote{假设在实际生产代码中，这些函数会基于……进行一些更高级的计算}。函数签名还使用了 \texttt{override} 关键字，这是 C++11 加入的。它告诉编译器某个成员函数即将被重写，因此编译器可以检查基类中是否存在对应的声明。关于这个关键字的更多信息，请参阅我的文章：Modern C++：用 override 和 final 实现安全与表现力 - C++ Stories\footnote{\url{https://www.cppstories.com/2021/override-final/}}。

现在我们可以编写同时使用这两种产品的代码：

\begin{codeboxt}[cpp]{代码清单 5.4 虚析构函数，主函数（在 Compiler Explorer 上运行）}
void CallCalculate(const Product& prod) {
std::cout << "calculating: " << prod.CalculateMass() << '\n';
}
int main() {
using std::unique_ptr;
using std::make_unique;
unique_ptr<Product> box = make_unique<BoxProduct>("box");
unique_ptr<Product> water = make_unique<FluidProduct>("water");
CallCalculate(*box.get());
CallCalculate(*water.get());
}
\end{codeboxt}

这个演示用例很简单，它创建了两个指向基类的智能指针，但它们被赋值为指向派生类的指针。当我们在 Compiler Explorer 上运行代码时，会看到以下输出：

\begin{codebox}[cpp]
box
water
calculating: 10
calculating: 100
water destructor...
box destructor...
\end{codebox}

如你所见，\texttt{CallCalculate(Product\& prod)} 函数接受一个指向基类的引用，然后可以调用它的函数。如果函数时虚函数，编译器会根据真实类型进行多态调用。

但是……你在这里看到错误了吗？

花点时间想一想……

看起来派生类的析构函数没有被调用！

这是因为我们用指针来持有对象。当智能指针离开作用域时，它会对指向基类的指针调用 \texttt{delete}。由于析构函数没有被标记为 \texttt{virtual}，多态不会生效，只会调用 \texttt{\textasciitilde Product()} 析构函数。

为了修复这一点，每个拥有虚函数的类也应该有一个虚析构函数：

\begin{codebox}[cpp]
virtual ~Product() {
std::cout << name_ << " destructor...\n";
}
\end{codebox}

这修复了我们的输出：

\begin{codebox}[cpp]
box
water
calculating: 10
calculating: 100
~FluidProduct...
water destructor...
~BoxProduct...
box destructor...
\end{codebox}

你可以在 Compiler Explorer 上试验这个正确的例子。

关于这个关键方面还有一条特定的 C++ Core 指南。参见 C.35：基类析构函数要么应该是公有的且虚的，要么是受保护的且非虚的：

\begin{myQuote}
理由：防止未定义行为。如果析构函数是公有的，那么调用代码可能试图通过基类指针销毁一个派生类对象，而如果基类的析构函数是非虚的，结果就是未定义的。
\end{myQuote}

\section{部分创建的对象}

编译器只为完全创建好的对象调用析构函数。请看对构造函数修改后的一版，它会检查 \texttt{id} 参数并抛出异常：

\begin{codebox}[cpp]
explicit Product(const char* name, unsigned id) : name_(name), id_(id) {
std::cout << name << ", id " << id << '\n';
if (id < 100)
throw std::runtime_error{"bad id..."};
}
\end{codebox}

\begin{codeboxt}[cpp]{代码清单 5.5 析构函数与部分对象创建（在 Compiler Explorer 上运行）}
int main() {
try {
Product tvset("TV Set", 123);
Product car("Mustang", 9);
}
catch (const std::exception& ex) {
std::cout << "exception: " << ex.what() << '\n';
}
}
\end{codeboxt}

当我们运行这个例子时，会得到以下输出：

\begin{codebox}[cpp]
TV Set, id 123
Mustang, id 9
TV Set destructor...
exception: bad id...
\end{codebox}

这一次例子创建了两个对象：\texttt{TV set} 和 \texttt{Mustang}。在输出中，我们可以注意到两个对象都调用了它们的构造函数，但只有一次析构函数调用（针对 \texttt{TV set}）。由于 \texttt{Mustang} 在构造函数中抛出了异常，析构函数不会被执行。

\begin{myWarnBox}
由于析构函数可能在编译器执行栈展开（stack unwinding）时被调用，它们不应该抛出异常，因为这可能导致调用 \texttt{std::terminate()}。阅读这条 C++ Core 指南建议以了解更多信息：E.16：析构函数、释放函数和 \texttt{swap} 绝不能失败。
\end{myWarnBox}

另一个重要方面是，要正确地管理在异常发生之前分配的资源。例如，如果你用原始指针动态分配了一些内存，可能会造成内存泄漏。看下面的示例：


\begin{codeboxt}[cpp]{代码清单 5.6 部分创建对象中的内存泄漏（在 Compiler Explorer 上运行）}
class Product {
public:
explicit Product(int id) : res_(new Resource()) {
std::cout << "Product: id " << id << '\n';
if (id < MIN_ID)
throw std::runtime_error{"bad id..."};
}
~Product() {
delete res_;
std::cout <<
"~Product...\n";
}
private:
Resource* res_;
};
int main() {
try {
Product invalid(MIN_ID - 1);
}
catch (const std::exception& ex) {
std::cout << "exception: " << ex.what() << '\n';
}
}
\end{codeboxt}

输出：

\begin{codebox}[cpp]
Product: id 99
exception: bad id...
\end{codebox}

如你所见，编译器没有为 \texttt{Resource} 调用析构函数，内存也没有被释放（因为我们在构造函数中调用了 \texttt{: res\_(new Resource())}）。我们可以在抛出之前 \texttt{delete res\_} 来修复这个泄漏。不过，手动资源管理很脆弱，最好寻找更好的解决方案。

修复这种资源泄漏的关键机制是依赖具有自动存储期（automatic storage duration）的变量和数据成员，比如普通的数值类型。在这种情况下，栈展开会正确地销毁它们并调用它们的析构函数。这就是为什么智能指针的析构函数可以被安全地调用：

\begin{codeboxt}[cpp]{代码清单 5.7 修复部分创建对象中的内存泄漏（在 Compiler Explorer 上运行）}
class Product {
public:
explicit Product(int id) : res_(std::make_unique<Resource>()) {
std::cout << "Product: id " << id << '\n';
if (id < MIN_ID)
throw std::runtime_error{"bad id..."};
}
~Product() {
std::cout <<
"~Product...\n";
}
private:
std::unique_ptr<Resource> res_; // << 现在是智能指针！
};
int main() {
try {
Product invalid(MIN_ID - 1);
}
catch (const std::exception& ex) {
std::cout << "exception: " << ex.what() << '\n';
}
}
\end{codeboxt}

当我们运行代码时，会看到以下输出：

\begin{codebox}[cpp]
Product: id 99
~Resource
exception: bad id...
\end{codebox}

编译器没有为 \texttt{Product} 调用析构函数，但栈展开正确地为所有具有自动存储期的数据成员调用了析构函数。

\begin{myInfoBox}
你可以在以下站点阅读更多关于栈展开和资源处理的信息：\texttt{throw} 表达式 - @C++Reference，以及 Exceptions and Error Handling, C++ FAQ。
\end{myInfoBox}


\section{编译器生成的析构函数}

与其他特殊成员函数一样，如果你不提供自己的实现，编译器会为你的类创建一个隐式的默认析构函数。基本规则是：每个数据成员以及基类都必须有一个可访问的析构函数（它们没有被删除、不可访问，也不模糊）。例如：

\begin{codeboxt}[cpp]{代码清单 5.8 编译器生成的析构函数（在 C++Insights 上运行）}
#include <iostream>
#include <string>
class Product {
public:
explicit Product(int id, const std::string& name)
: id_(id), name_(name) {
std::cout << "Product(): " << id_ << ", " << name_ << '\n';
}
private:
int id_;
std::string name_;
};
int main() {
Product first{10, "basic"};
}
\end{codeboxt}

在 C++Insights 上，我们可以看到编译器的输出，以及它如何“看待”这段代码。如你所注意到的，编译器为我们创建了以下析构函数：

\begin{codebox}[cpp]
inline ~Product() noexcept = default;
\end{codebox}

你可以在 C++ 参考文档的“隐式声明的析构函数”一节阅读更多关于编译器生成的析构函数的内容。

\section{总结与用例}

本章介绍了析构函数，这是对象结束生存期时被调用的特殊成员函数。在大多数情况下，我们可以用这个能力来正确地清理对象，并释放我们已使用但尚未释放的任何资源。

例如，你在对象创建时分配了一些内存，然后必须释放该内存以避免内存泄漏。类似地，你可以打开一个文件或数据库连接，然后必须确保在对象离开作用域时关闭该文件或连接。

析构函数是 C++ 语言中最好的特性之一，因为它提供了一个明确且定义良好的调用点。这与动态垃圾回收相反，后者可能在后台运行，可能拖慢程序，并且对过程的控制更少。析构函数也是 C++ 中一个流行术语的关键要素，称为 RAII（资源获取即初始化，Resource Acquisition is Initialization），由 C++ 语言的作者 B. Stroustrup 提出。它指出，持有资源是类的不变式，并与对象的生存期绑定在一起。更多信息请参阅 Wikipedia\footnote{\url{https://en.wikipedia.org/wiki/Resource_acquisition_is_initialization}}。

幸运的是，在 Modern C++ 中，需要自定义析构函数的地方越来越少。例如，当你的数据成员是标准容器（如 \texttt{std::vector<int>}，或 \texttt{std::map<std::string, int>}）时，你可以依赖默认析构函数来完成工作。像 \texttt{std::vector<int>} 这样的标准容器可能会分配内存缓冲区，但它们也会管理该缓冲区并正确释放它，所以你在类中使用它们时不需要采取任何操作。
